

%%% fhc mu-like

\begin{table}[htb]
\centering
\begin{tabular}{| c | c | c |}
\hline
  {\bf Source of uncertainty}    & \multicolumn{1}{c|}{{\bf ${\delta}N_{SK}$/$N_{SK}$}}\\
                                 & \multicolumn{1}{c|}{(Asimov B)}\\

\hline
\hline  SKDet+FSI+SI                   & 4.02\%                        \\
\hline
  Flux                           & 3.56\%                        \\
\hline
\hline
  2p-2h (corr)                     & 3.49\%                        \\
  2p-2h-bar (corr)                  & 0.20\%                        \\
  NC other (uncorr)              & 0.77\%                        \\
  NC 1gamma (uncorr)             & 0.00\%                        \\
  XSec nue/numu (uncorr)                 & 0.01\%                        \\
\hline  
  XSec Tot (corr)                  & 4.03\%                        \\
  XSec Tot                       & 4.10\%                        \\
\hline
\hline
  Flux+XSec (All)                      & 2.85\%                        \\
\hline
  Flux+XSec+SKDet+FSI+SI         & 4.96\%                        \\
\hline
\hline
  Oscillations                   & 0.10\%                        \\
\hline
\hline
  All                            & 4.96\%                        \\
\hline
\hline

\end{tabular}
\caption{
 Effect of 1$\sigma$ variations of the groups of systematics parameters on total number of
 1-ring muon-like events for neutrino mode obtained by performing 10k toys MC.
The systematic parameter values measured with the ND280 data 
shown in 
\cref{tab:banff_uncertainties_flux_fhc},
~\ref{tab:banff_uncertainties_flux_rhc}
~\ref{tab:banff_uncertainties_xsec},
the nominal Super-K systematic parameters shown in \cref{tab:banff_uncertainties_skdetfsi}
and the oscillation parameters of the Asimov data set B shown in \cref{tab:osc_params_joint_asimov}
are used.
Since the parameters $\delta_{CP}$, $\theta_{23}$ and $\Delta m^2_{32}$ ($\Delta m^2_{13}$) and mass hierarchy
will be measured,
only the uncertainty on the parameters 
$\sin^2 \theta_{13}$,
$\sin^2 \theta_{12}$
and 
$\Delta m^2_{21}$ 
are taken into account.
}                                                                                                                                                                                          
\label{tab:nsk_systematic_table_summary_joint_asimov_B_fhc_mulike}
\end{table}




%%%%% fhc e-like


\begin{table}[htb]
\centering
\begin{tabular}{| c | c | c |}
\hline
  {\bf Source of uncertainty}    & \multicolumn{1}{c|}{{\bf ${\delta}N_{SK}$/$N_{SK}$}}\\
                                 & \multicolumn{1}{c|}{(Asimov B)}\\

\hline
\hline  SKDet+FSI+SI                   & 3.67\%                        \\
\hline
  Flux                           & 3.61\%                        \\
\hline
\hline
  2p-2h (corr)                     & 3.94\%                        \\
  2p-2h-bar (corr)                  & 0.07\%                        \\
  NC other (uncorr)              & 0.21\%                        \\
  NC 1gamma (uncorr)             & 1.94\%                        \\
  XSec nue/numu (uncorr)                 & 2.55\%                        \\
\hline
  XSec Tot (corr)                  & 4.16\%                        \\
  XSec Tot                       & 5.26\%                        \\
\hline
\hline
  Flux+XSec                      & 4.29\%                        \\
\hline
\hline
  Flux+XSec+SKDet+FSI+SI         & 5.68\%                        \\
\hline
\hline
  Oscillations                   & 4.18\%                        \\
\hline
\hline
  All                            & 7.19\%                        \\
\hline
\hline

\end{tabular}
\caption{
 Effect of 1$\sigma$ variations of the groups of systematics parameters on total number of
 1-ring electron-like events for neutrino mode obtained by performing 10k toys MC.
The systematic parameter values measured with the ND280 data 
shown in 
\cref{tab:banff_uncertainties_flux_fhc},
~\ref{tab:banff_uncertainties_flux_rhc}
~\ref{tab:banff_uncertainties_xsec},
the nominal Super-K systematic parameters shown in \cref{tab:banff_uncertainties_skdetfsi}
and the oscillation parameters of the Asimov data set B shown in \cref{tab:osc_params_joint_asimov}
are used.
Since the parameters $\delta_{CP}$, $\theta_{23}$ and $\Delta m^2_{32}$ ($\Delta m^2_{13}$) and mass hierarchy
will be measured,
only the uncertainty on the parameters 
$\sin^2 \theta_{13}$,
$\sin^2 \theta_{12}$
and 
$\Delta m^2_{21}$ 
are taken into account.
}                                                                                                                                                                                          
\label{tab:nsk_systematic_table_summary_joint_asimov_B_fhc_elike}
\end{table}



%%%%%% rhc mu-like


\begin{table}[htb]
\centering
\begin{tabular}{| c | c | c |}
\hline
  {\bf Source of uncertainty}    & \multicolumn{1}{c|}{{\bf ${\delta}N_{SK}$/$N_{SK}$}}\\
                                 & \multicolumn{1}{c|}{(Asimov B}\\

\hline
\hline  SKDet+FSI+SI                   & 3.83\%                        \\
\hline
  Flux                           & 3.70\%                        \\
\hline
\hline
  2p-2h (corr)                     & 2.95\%                        \\
  2p-2h-bar (corr)                  & 1.83\%                        \\
  NC other (uncorr)              & 0.73\%                        \\
  NC 1gamma (uncorr)             & 0.00\%                        \\
  XSec nue/numu (uncorr)                  & 0.00\%                        \\
\hline
  XSec Tot (corr)                  & 4.10\%                        \\
  XSec Tot                       & 4.17\%                        \\
\hline
\hline
  Flux+XSec                      & 3.31\%                        \\
\hline
\hline
  Flux+XSec+SKDet+FSI+SI         & 5.04\%                        \\
\hline
\hline
  Oscillations                   & 0.06\%                        \\
\hline
\hline
  All                            & 5.11\%                        \\
\hline
\hline

\end{tabular}
\caption{
 Effect of 1$\sigma$ variations of the groups of systematics parameters on total number of
 1-ring muon-like events for antineutrino mode obtained by performing 10k toys MC.
The systematic parameter values measured with the ND280 data 
shown in 
\cref{tab:banff_uncertainties_flux_fhc},
~\ref{tab:banff_uncertainties_flux_rhc}
~\ref{tab:banff_uncertainties_xsec},
the nominal Super-K systematic parameters shown in \cref{tab:banff_uncertainties_skdetfsi}
and the oscillation parameters of the Asimov data set B shown in \cref{tab:osc_params_joint_asimov}
are used.
Since the parameters $\delta_{CP}$, $\theta_{23}$ and $\Delta m^2_{32}$ ($\Delta m^2_{13}$) and mass hierarchy
will be measured,
only the uncertainty on the parameters 
$\sin^2 \theta_{13}$,
$\sin^2 \theta_{12}$
and 
$\Delta m^2_{21}$ 
are taken into account.
}                                                                                                                                                                                          
\label{tab:nsk_systematic_table_summary_joint_asimov_B_rhc_mulike}
\end{table}



%%%%%% rhc e-like

\begin{table}[htb]
\centering
\begin{tabular}{| c | c | c |}
\hline
  {\bf Source of uncertainty}    & \multicolumn{1}{c|}{{\bf ${\delta}N_{SK}$/$N_{SK}$}}\\
                                 & \multicolumn{1}{c|}{(Asimov B)}\\

\hline
\hline  SKDet+FSI+SI                   & 3.88\%                        \\
\hline
  Flux                           & 3.81\%                        \\
\hline
\hline
  2p-2h (corr)                     & 2.92\%                        \\
  2p-2h-bar (corr)                  & 2.56\%                        \\
  NC other (uncorr)              & 0.32\%                        \\
  NC 1gamma (uncorr)             & 2.81\%                        \\
  XSec nue/numu (uncorr)                 & 1.68\%                        \\
\hline  
  XSec Tot (corr)                  & 4.52\%                        \\
  XSec Tot                       & 5.59\%                        \\
\hline
\hline
  Flux+XSec                      & 4.62\%                        \\
\hline
\hline
  Flux+XSec+SKDet+FSI+SI         & 6.18\%                        \\
\hline
\hline
  Oscillations                   & 3.59\%                        \\
\hline
\hline
  All                            & 7.09\%                        \\
\hline
\hline

\end{tabular}
\caption{
 Effect of 1$\sigma$ variations of the groups of systematics parameters on total number of
 1-ring electron-like events for antineutrino mode obtained by performing 10k toys MC.
The systematic parameter values measured with the ND280 data 
shown in 
\cref{tab:banff_uncertainties_flux_fhc},
~\ref{tab:banff_uncertainties_flux_rhc}
~\ref{tab:banff_uncertainties_xsec},
the nominal Super-K systematic parameters shown in \cref{tab:banff_uncertainties_skdetfsi}
and the oscillation parameters of the Asimov data set B shown in \cref{tab:osc_params_joint_asimov}
are used.
Since the parameters $\delta_{CP}$, $\theta_{23}$ and $\Delta m^2_{32}$ ($\Delta m^2_{13}$) and mass hierarchy
will be measured,
only the uncertainty on the parameters 
$\sin^2 \theta_{13}$,
$\sin^2 \theta_{12}$
and 
$\Delta m^2_{21}$ 
are taken into account.
}                                                                                                                                                                                          
\label{tab:nsk_systematic_table_summary_joint_asimov_B_rhc_elike}
\end{table}


















\begin{table}
\center
\begin{tabular}{| c || c | c || c | c |}
\hline
Source of uncertainty 
&$\mu$-like ${\delta} \Big( \frac{\# {\nu} \text{-mode}}{\# {\bar{\nu}} \text{-mode} } \Big) $ / $\Big \langle \frac{\# {\nu} \text{-mode}}{\# {\bar{\nu}} \text{-mode} } \Big \rangle$ 
&$e$-like ${\delta} \Big( \frac{\# {\nu} \text{-mode}}{\# {\bar{\nu}} \text{-mode} } \Big) $ / $\Big \langle \frac{\# {\nu} \text{-mode}}{\# {\bar{\nu}} \text{-mode} } \Big \rangle$ \\

\hline
SKDet & 0.6\% & 1.2\% \\ \hline
FSI+SI & 2.5\%  & 3.6\% \\ \hline
SKDet+FSI+SI & 2.6\% & 3.7\% \\ \hline
\hline
Flux & 1.6\%  & 1.7\% \\ \hline
XSec NC other (uncorr)  & 0.0\%  & 0.1\% \\ \hline
XSec NC $1\gamma$ (uncorr)  & 0.0\%  & 0.9\% \\ \hline
XSec $\nu_e$ / $\nu_{\mu}$ (uncorr)  & 0.0\%  & 3.3\% \\ \hline
XSec Tot (All) & 1.7\%  & 4.4\% \\ \hline
 \hline
Flux+XSec  & 1.8\%  & 4.1\% \\ \hline
 \hline
Flux+XSec+SDet+FSI+SI  & 3.1\%  & 5.6\% \\ \hline
 \hline
Oscillations  & 0.0\%  & 0.6\% \\ \hline
 \hline
All  & 3.2\%  & 5.6\% \\ \hline

\end{tabular}
\caption{
 Effect of 1$\sigma$ variations of the several source of systematic uncertainties 
 on the ratio between the
 total number of events in neutrino and antineutrino modes
 for 1-ring electron- and muon-like samples
 obtained by performing 10k toys MC.
 The systematic parameter values measured with the ND280 data 
shown in 
\cref{tab:banff_uncertainties_flux_fhc},
~\ref{tab:banff_uncertainties_flux_rhc}
~\ref{tab:banff_uncertainties_xsec},
the nominal Super-K systematic parameters shown in \cref{tab:banff_uncertainties_skdetfsi}
and the oscillation parameters of the Asimov data set B shown in \cref{tab:osc_params_joint_asimov}
are used.
Since the parameters $\delta_{CP}$, $\theta_{23}$ and $\Delta m^2_{32}$ ($\Delta m^2_{13}$) and mass hierarchy
will be measured,
only the uncertainty on the parameters 
$\sin^2 \theta_{13}$,
$\sin^2 \theta_{12}$
and 
$\Delta m^2_{21}$ 
are taken into account.
}             
\label{tab:ratio_systematic_table_summary_asimov_B}                                                                                                                                                                             
\end{table}
